\documentclass[a4paper,12pt]{article}
\setlength{\oddsidemargin}{-0.33in}
\setlength{\evensidemargin}{0in}
\setlength{\textwidth}{7in}
\setlength{\textheight}{10.2in}
\setlength{\footskip}{.5in}

\setlength{\voffset}{-.5in}
\setlength{\topmargin}{-0.5in}
\setlength{\headheight}{.5in}
\setlength{\headsep}{12pt}

\setlength\parindent{0pt}

% \usepackage{savetrees}
\usepackage{stmaryrd}
\usepackage{hyperref}
\usepackage{color}
\usepackage{outlines}
\usepackage{rotating}
\usepackage{overpic}
\usepackage{algorithm}
\usepackage{algorithmic}
\usepackage{amsmath}
\usepackage{amssymb}
\usepackage{fancyhdr,lastpage}
\usepackage{enumerate}
\usepackage{graphicx}
\usepackage{epstopdf}
\usepackage{rotating}
\usepackage{mathrsfs}
\usepackage{tipa}
\usepackage{setspace}
\usepackage{natbib}
\usepackage{cancel}


%\doublespacing
\newcommand{\bs}[1]{\boldsymbol{#1}}
\newcommand{\real}[1]{\Re\left\{#1\right\}}
\newcommand{\TD}[2]{\frac{D#1}{D#2}}
\newcommand{\td}[2]{\frac{d#1}{d #2}}
\newcommand{\pd}[2]{\frac{\partial #1}{\partial #2}}
\newcommand{\pdd}[3]{\frac{\partial^2 #1}{\partial #2 \partial #3}}
\newcommand{\pdt}[2]{\frac{\partial^2 #1}{\partial #2 ^2}}
\newcommand{\expv}[1]{\left \langle #1 \right \rangle}
\newcommand{\tensort}[1]{\underline{\underline{#1}}}
\newcommand{\erf}[1]{\mbox{erf}{\left( {#1}\right)}}
\newcommand{\erfc}[1]{\mbox{erfc}{\left( {#1}\right)}}
\newcommand{\ands}{ \ \ \ \mbox{ and } \ \ \ }
\newcommand{\spword}[1]{ \ \ \ \mbox{ {#1} } \ \ \ }

\begin{document}
\title{Polythermal Equation Solution methods}
\author{Toby Harvey}
\maketitle
\vspace{-1cm}

\section{Equation Derivations}

Assume that the temperate ice is a mixture of ice with density $\rho_i$ and velocity $u$, and water with volume fraction $\phi$, density $\rho_w$ and velocity $u_w = u + u_l$ where $u_l$ is the additional velocity of the water beyond the ice. Lastly assume that ice will melt at a ratet $m$ and turn directly into water. Then, conservation of mass for both ice and water give:

\begin{align*}
  &\frac{D(\rho_i(1 - \phi))}{Dt} = -m & (\text{ice mass balanace})\\
  &\rho_i\frac{D((1 - \phi))}{Dt} = -m \\
  &\pd{(1 - \phi)}{t} + \nabla \cdot[(1-\phi)u] = -\frac{m}{\rho_i} \tag{1} \label{eq:ice-mass-balance}\\
  &\frac{D\rho_w\phi}{Dt} = m & (\text{water mass balance})\\
  &\pd{\phi}{t} + \nabla \cdot [\phi(u + u_l)] = \frac{m}{\rho_w} & (\text{define } j = \phi u_l)\\
  &\pd{\phi}{t} + \nabla \cdot [\phi u + j] = \frac{m}{\rho_w} \tag{2} \label{eq:water-mass-balance} \\
\end{align*}

Because there is melting latent heat $L$ is being absorbed, and similar to cold ice energy is being produced by viscous dissipation $a$, because the temperature in the temperate region is by defintion $T_m$, a constant, overall conservation of energy gives:

\begin{align*}
  \cancel{\text{Temperature terms}}= a - Lm \rightarrow m = \frac{a}{L} \tag{3} \label{eq:conservation-of-energy}
\end{align*}

The flux of liquid water $j$ follows darcy's law (assume the viscous resistive force is linear to velocity in stokes equation):

\begin{align*}
  j = \Pi(\phi)(\rho_w g - \nabla p_w) =  \kappa\phi^\alpha(\rho_w g - \nabla p_w)\tag{4} \label{eq:darcys-law}
\end{align*}
Where the hydraulic conductivity follows $\Pi(\phi) =\kappa \phi^\alpha$, a Kozeny–Carman relation, and $\kappa$ is the conductivity.

Since the ice matrix is no-longer incompressable it exerts a compaction pressure $p_e$, where the total pressure is $p = p_e + p_w$, and:

\begin{align*}
  p_e = - \zeta (\nabla \cdot u) \tag{5} \label{eq:compaction-pressure}
\end{align*}

where $\zeta = \frac{\eta}{\phi}$ is the bulk viscosity of ice, and $\eta$ is the viscosity of ice.


using mass conversation, energy conservation, darcy's law and the compaction pressure relation, we can derive the equations below. Adding \eqref{eq:ice-mass-balance}, to \eqref{eq:water-mass-balance} gives:

\begin{align*}
  &\nabla \cdot (u - \phi u) + \nabla \cdot(\phi u + j) = \frac{m}{\rho_w} - \frac{m}{\rho_i}\\
  &\nabla \cdot u + \nabla \cdot j = \frac{\rho_i - \rho_w}{\rho_w\rho_i}m \tag{6} \label{eq:sum-mass-conservation}\\
\end{align*}

Now taking a boussinesq approximation so that $\rho_w = \rho_i$, \eqref{eq:sum-mass-conservation} gives:

\begin{align*}
  \nabla \cdot j = - \nabla \cdot u
\end{align*}

plugging in \eqref{eq:compaction-pressure} gives:

\begin{align*}
  \nabla \cdot j = \frac{p_e\phi}{\eta}
\end{align*}

Now plugging this into \eqref{eq:water-mass-balance}, and using \eqref{eq:conservation-of-energy} gives:

\begin{align*}
  &\pd{\phi}{t} + \nabla \cdot [\phi u] + \frac{p_e\phi}{\eta} = \frac{a}{\rho L} \\
\end{align*}

Plugging this same relationship into darcy's law and replacing $p_w$ with $p - p_e$ gives:

\begin{align*}
  &\frac{p_e\phi}{\eta} = \nabla \cdot (\kappa\phi^\alpha(\rho_w g - \nabla (p - p_e))) & (\text{substituting } \nabla p = \rho g + \nabla p_r)\\
  &\frac{p_e\phi}{\eta} = \nabla \cdot (\kappa\phi^\alpha((\rho_w - \rho) g + \nabla p_r + \nabla p_e))\\
\end{align*}

where from scaling arguments $\nabla p_r$ can be ignored. All together the temperate equations are then:

\begin{align*}
  &\frac{p_e\phi}{\eta} = \nabla \cdot (\kappa\phi^\alpha((\rho_w - \rho) g + \nabla p_e))\\
  &\pd{\phi}{t} + \nabla \cdot [\phi u] + \frac{p_e\phi}{\eta} = \frac{a}{\rho L} \\
\end{align*}

below are the one versions with constant velocity.

\section{1D - 3 equation system}

\subsection{Cold Ice}

\subsubsection{Temperature Equation}
\begin{align*}
  \rho c \left(\pd{T}{t} + u \pd{T}{z}\right) - \pd{}z\left(k \pd{T}{z}\right) = a \\
  T(S,t) = T_0\\
  \pd{T(\Gamma,t)}{z} = T(\Gamma,t) = 0
\end{align*}

\begin{center}
\begin{tabular}{ |c|c|c| } 
 \hline
  Parameter & Name & units\\
  \hline
  $\rho$ & ice density & $\frac{kg}{m^3}$\\
  \hline
  $c$ & heat capacity of ice & $\frac{J}{kg K^\circ}$\\
  \hline
  $K$ & Thermal conductivity & $\frac{J}{ ms K^\circ}$\\
  \hline
  $a$ & viscous dissipation & $\frac{J}{m^3 s}$\\
  \hline
  $T_0$ & surface temperature & $K^\circ$\\
  \hline
  $\Gamma$ & cold-temperate interface & \\
  \hline
\end{tabular}
\end{center}

\subsection{Temperate Ice}

\subsubsection{Compaction Pressure ($P_e$) equation}

\begin{align*}
  &\pd{}{z}\left(\kappa\phi^\alpha\left((\rho_w - \rho)g + \pd{}{z}p_e\right)\right) = \frac{\phi p_e}{\eta}\\
  &p_e(0,t) = N_0\\
  &\left(\kappa \phi^\alpha\left( (\rho_w - \rho)g + \pd{}{z} p_e\right)\right)\bigg\rvert_{z=\Gamma, t=t} = 0\\
\end{align*}

\subsubsection{Porosity ($\phi$) equation}

\begin{align*}
  &\pd{\phi}{t} + u\pd{\phi}{z} + \frac{\phi p_e}{\eta} = \frac{a}{\rho L}\\
  &\phi(\Gamma, t) = 0\\
\end{align*}


\begin{center}
\begin{tabular}{ |c|c|c| } 
 \hline
  Parameter & Name & units\\
  \hline
  $\kappa$ & permeability & $\frac{m^2}{Pa s}$\\
  $\alpha$ & Carman–Kozeny coefficent & unitless\\
  $\eta$ & viscosity & $Pa s$\\
  $L$ & Latent heat of freezing & $\frac{J}{kg}$\\
  \hline
\end{tabular}
\end{center}

\section{1D - ethalpy system}

Define enthalpy as:

\begin{align*}
  h = \rho c T + \rho L \phi
\end{align*}

Then the temperature equation and porosity equation combine to:

\begin{align*}
\pd{h}{t} + u \pd{h}{z} - \pd{}{z}\left(\frac{k}{\rho c}\pd{T}{z}\right) + \frac{\rho L}{\eta}\phi p_e = a
\end{align*}

Since diffusion only occurs in the cold domain, and porosity only effects the temperate domain this is better written as:

\begin{align*}
  &h = \rho c T + \rho L \phi\\
  \\
  &\text{Enthalpy Equation:}\\
  &\pd{h}{t} + u \pd{h}{z} + Q(h) = a\\
  &Q =
  \begin{cases}
    - \pd{}{z}\left(k\pd{h}{z}\right), & h < h_{melt}\\
    \frac{hp_e}{\eta}, & h \geq h_{melt}
  \end{cases}\\
  \\
  &\text{Initial and Boundary Conditions:}\\
  &h(z, 0) = I(z)\\
  &h(S, t) = T_{surf}\\
  &h(\Gamma(t), t) = 0\\
  \\
  &\text{Compaction Pressure Equation:}\\
  &\pd{}{z}\left(\kappa\phi^\alpha\left((\rho_w - \rho)g + \pd{}{z}p_e\right)\right) = \frac{\phi p_e}{\eta}\\
  &\text{Boundary Conditions:}\\
  &p_e(0,t) = N_0\\
  &\left(\kappa \phi^\alpha\left( (\rho_w - \rho)g + \pd{}{z} p_e\right)\right)\bigg\rvert_{z=\Gamma, t=t} = 0\\
\end{align*}


\section{Solution Methods}

\subsection{Spatial Discretization}

using a standard FEM discretization the enthalpy operators are:

\begin{align*}
  &\text{Enthalpy Operators}:\\
  \\
  &\pd{h}{t} \approx \sum_{j=1}^N \int_{\Omega} \psi_i\psi_j dz\pd{\hat{h}_j}{t} = M\pd{\hat{h}}{t}, ~~~ i=1,...N\\
  \\
  & u\pd{h}{z} \approx \sum_{j=1}^N\int_{\Omega_T} \psi_iu\pd{\psi_j}{z}dz \hat{h}_j = S\hat{h} , ~~~ i=1,...N\\
  \\
  &-\pd{}{z}\left(k\pd{h}{z}\right) \underbrace{\approx}_{\text{IVP gives + sign}} \sum_{j=1}^N\int_{\Omega} \pd{\psi_i}{x}k\pd{\psi_j}{x} dz \hat{h}_j = K_k\hat{h}, ~~~ i = \Gamma_n(\hat{h}),...N\\
  \\
  &\frac{hp_e}{\eta} \approx \int_{\Omega} \psi_i \frac{p_e}{\eta}\psi_j dz \hat{h}_j = M_{p_e}\hat{h}, ~~~ i = 1,...\Gamma_n(\hat{h}) \\
  \\
  &Q(\hat{h})\hat{h} = \begin{bmatrix}
    M_{p_e} & 0\\
    0    & K_k
    \end{bmatrix} \hat{h}
  \\
  \\
  & a \approx \int_{\Omega} a\psi_i dz = F_a, ~~~ i = 1,...N \\
  \\
  &\text{Compaction Pressure Operators}:\\
  \\
  &\pd{}{z}(\kappa\phi^\alpha(\rho_w - \rho)g) = \kappa(\rho_w - \rho)g\pd{\phi^\alpha}{z} \approx \kappa(\rho_w - \rho)g \int_\Omega\phi^\alpha\pd{\psi_i}{z} = \kappa\Delta\rho g F_{\phi^\alpha}, ~~~ i=1,...\Gamma_n(\hat{h})\\
  \\
  & \kappa\pd{}{z}\left(\phi^\alpha\pd{p_e}{z}\right) \underbrace{\approx}_{\text{IVP gives - sign}} -\sum_{j=1}^N\int_{\Omega} \pd{\psi_i}{x}\kappa(\hat{\phi}_j^\alpha\psi_j)\pd{\psi_j}{x} dz \hat{p}_{e,j} = -\kappa K_{\phi^\alpha}\hat{p}_e , ~~~ i=1,...\Gamma_n(\hat{h})\\
  \\
  &\frac{\phi p_e}{\eta} \approx \sum_{j=1}^N\int_{\Omega}\psi_i \frac{(\hat{\phi}_j\psi_j)}{\eta}\psi_j dz \hat{p}_{e,j} = M_{\phi} \hat{p_e} , ~~~ i=1,...\Gamma_n(\hat{h})
\end{align*}

Where $\Gamma_n$ is the node where the cold/temperate boundary is. With these operators the enthalpy equation becomes:

\begin{align*}
  M\pd{\hat{h}}{t} + S\hat{h} + Q(\hat{h})\hat{h} - F_a = 0
\end{align*}

and the compaction pressure equation becomes:

\begin{align*}
  -(K_{\phi^\alpha} + M_\phi)\hat{p}_e = \kappa\Delta\rho g F_{\phi^\alpha}
\end{align*}

$p_e$ is a linear solve assuming $\Gamma(\hat{h})$ is known.

\subsection{Timestepping for Enthalpy}

\subsubsection{Explict methods}

Since we have diffusion we are at least limited (assuming forward euler and second order central difference - not what we are doing but maybe representative) (Von Neumann analysis) to: $\Delta  t \leq \frac{\Delta x^2}{2k}$, this just gives:

\begin{align*}
  \hat{h}^{n+1} = M^{-1}\mathcal{T}(\hat{h}^n)
\end{align*}

where $\mathcal{T}$ is the operator for any explict timestepping scheme, and $M^{-1}$ is the inverse of a lumped (diagonal) mass matrix so we don't need to solve a linear system. Two major disadvantages to this: stiffness with $\Delta x^2$ scaling, and no treatment of the non-linear $Q(\hat{h})\hat{h}$ term.

\subsection{Operator splitting method}

Taking just the diffusive part first and solve over the entire domain with crank-nicolson noting that $T = T_{melt}$ in the temperate domain for a timestep of $\frac{\Delta t}{2}$:

\begin{align*}
  &M\frac{2}{\Delta t}\left(\hat{T}^{n + \frac{1}{3}} - \hat{T}^n\right) = \frac{1}{2}\left(-K_k\hat{T}^{n + \frac{1}{3}} - K_k\hat{T}^{n}\right)\\
  &\hat{T}^{n + \frac{1}{3}} = \left(M + \frac{\Delta t}{4}K_k\right)^{-1}\left(m - \frac{\Delta t}{4}K_k\right) \hat{T}^n\\
\end{align*}

Then reset enthalpy to be $\hat{h}^{n+\frac{1}{3}} = \hat{\phi}^{n} + \hat{T}^{n+\frac{1}{3}}$, then solve the advective and porosity part explictly for $\Delta T$:

\begin{align*}
  \hat{h}^{n + \frac{2}{3}} = M_{lump}^{-1}\mathcal{T}(\hat{h}^{n + \frac{1}{3}}, \hat{p}_e^n, \Delta t,S, M_{\phi})
\end{align*}

Lastly do the source term:

\begin{align*}
  \hat{h}^{n + 1} = M_{lump}^{-1} \mathcal{T}(\hat{h}^{n + \frac{2}{3}}, F_a)
\end{align*}

Error analysis for this?


\subsection{Picard iterations}

Solve the entire enthalpy system implicitly (Crank–Nicolson), and picard iterate on $\hat{h}$, and $p_e$ so the non-linearity in $Q$ is resolved:

\begin{align*}
  &\hat{h}^{n+1, k+1} = \left(M + \frac{\Delta t}{2} \left( S + Q(\hat{h}^{n+1, k})\right)\right)^{-1}\left(M - \frac{\Delta t}{2}\left(S + Q(\hat{h}^n\right)\right)\hat{h}^n + \Delta t F_a\\
  &\hat{p}_e^{k+1} = -((K_{\phi^\alpha} + M_\phi)^k)^{-1}(\kappa\Delta\rho g) F_{\phi^\alpha}^k
\end{align*}

Where we have to couple the compaction pressure to the iteration as well, because it influances $Q(\hat{h})$.



\end{document}


